\rok{Децембар 2025.}{D}

Први практични тест из Алгоритама и структура података

\zadatak{1}{Selection Sort}
Написати Selection Sort алгоритам за сортирање целобројног низа дужине $n$.

\textbf{Додатне напомене за решавање задатка:}
\begin{itemize}
	\item Улаз је несортиран низ целих бројева.
	\item Излаз је сортирани низ целих бројева.
	\item У template-у се налази класа \texttt{RandomArrayGenerator} коју треба искористити за генерисање низа случајних бројева.
	\item Креирати класу \texttt{Selection} и у оквиру ње генерисати методу:
	\begin{Verbatim}[fontsize=\small]
public static void Sort(long[] arr)
	\end{Verbatim}
\end{itemize}

\zadatak{2}{Минималан збир близанаца}
У повезаној листи са парним бројем чворова величине $n$, дефинишемо близаначки пар чворова на следећи начин:

За $i$-ти чвор (почевши од индекса 0) постоји одговарајући близанац на позицији $(n-1-i)$, ако важи $0 <= i <= (n/2)-1$.

На пример:
\begin{itemize}
	\item Када је $n=4$, чвор са индексом 0 има близанца у чвору са индексом 3, док чвор са индексом 1 има близанца у чвору са индексом 2. Ова два пара су једини близаначки парови за листу величине 4.
\end{itemize}

Збир близанаца дефинише се као збир вредности чвора и његовог близанца. Задатак је да пронађете најмањи збир близаначких парова у листи. Решење треба да пронађе минималан збир близанаца у $O(n)$ времену.

\textbf{Додатни захтеви за решавање задатка:}
\begin{itemize}
	\item Алгоритам писати у оквиру класе \texttt{LinkedList}.
	\item Захтеван потпис методе:
	\begin{Verbatim}[fontsize=\small]
public int findMinTwinSum()
	\end{Verbatim}
\end{itemize}

\textbf{Пример 1:}
\begin{itemize}
	\item \textbf{Улаз:} $5 \leftrightarrow 4 \leftrightarrow 1 \leftrightarrow 1$ \\
	      \textbf{Излаз:} $5$
\end{itemize}

\textbf{Пример 2:}
\begin{itemize}
	\item \textbf{Улаз:} $4 \leftrightarrow 2 \leftrightarrow 2 \leftrightarrow 3$ \\
	      \textbf{Излаз:} $4$
\end{itemize}

\zadatak{3}{Asteroid Collision}
Налазите се у свемирској станици и пратите низ астероида који се крећу у истој линији. Сваки астероид има своју величину (апсолутна вредност броја) и смер кретања (позитиван број значи кретање десно, негативан лево). Када се два астероида сударе (крећу се један ка другом), мањи експлодира и нестаје. Ако су исте величине, оба нестају. Астероиди који се крећу у истом смеру се никада не сударају.

\textbf{Проблем:} Дат је низ целих бројева који представља астероиде. Потребно је симулирати све сударе и вратити стање низа астероида који су "преживели" након што се све колизије заврше.

\textbf{Додатни захтеви за решавање задатка:}
\begin{enumerate}
	\item Задатак решавати у оквиру методе:
	\begin{Verbatim}[fontsize=\small]
public int[] SimulateCollisions(int[] asteroids)
	\end{Verbatim}
	\item Користити Stack структуру података за симулацију тока судара:
	\begin{itemize}
		\item Итерирате кроз астероиде.
		\item Нови астероид се "бори" са онима на врху стека све док не буде уништен или док не уништи све којима прети судар.
	\end{itemize}
	\item Задатак решити са временском комплексношћу $O(n)$.
	\item Резултат треба да буде низ преосталих астероида у оригиналном редоследу кретања (с лева на десно).
\end{enumerate}

\textbf{Пример:}

\textbf{Улаз:} \texttt{[5, 10, -5, 2, -12, 8]}

\textbf{Кораци:}
\begin{itemize}
	\item \textbf{Корак 1.} Астероид: \texttt{5}; акција: \texttt{Push}; стање стека (од дна ка врху): \texttt{[5]}
	\item \textbf{Корак 2.} Астероид: \texttt{10}; акција: \texttt{Push}; стање стека (од дна ка врху): \texttt{[5, 10]}
	\item \textbf{Корак 3.} Астероид: \texttt{-5}; акција: \texttt{Судар}; стање стека (од дна ка врху): \texttt{[5, 10]}
	\item \textbf{Корак 4.} Астероид: \texttt{2}; акција: \texttt{Push}; стање стека (од дна ка врху): \texttt{[5, 10, 2]}
	\item \textbf{Корак 5.} Астероид: \texttt{-12}; акција: \texttt{Судар}; стање стека (од дна ка врху): \texttt{[5, 10]}
	\item \textbf{Корак 5а.} Астероид: \texttt{-12}; акција: \texttt{Судар}; стање стека (од дна ка врху): \texttt{[5]}
	\item \textbf{Корак 5б.} Астероид: \texttt{-12}; акција: \texttt{Судар}; стање стека (од дна ка врху): \texttt{[]}
	\item \textbf{Корак 5ц.} Астероид: \texttt{-12}; акција: \texttt{Push}; стање стека (од дна ка врху): \texttt{[-12]}
	\item \textbf{Корак 6.} Астероид: \texttt{8}; акција: \texttt{Push}; стање стека (од дна ка врху): \texttt{[-12, 8]}
\end{itemize}

\textbf{Излаз:} \texttt{[-12, 8]}

\sablon{Шаблон првог теста децембар 2025. група D}{2025/Decembar/GrupaD/AISP2025_Test1_GrupaD.linq}
